% =====================================================================
%  CONCLUSIONES Y RECOMENDACIONES
% =====================================================================
\chapter*{Conclusiones y recomendaciones}
\addcontentsline{toc}{chapter}{CONCLUSIONES Y RECOMENDACIONES}
\setcounter{section}{0}
\renewcommand{\thesection}{\arabic{section}}

\guia{Las conclusiones vinculan el desarrollo del trabajo con la pregunta del problema. Debe haber una conclusión por cada objetivo específico, y la conclusión general debe decir si se acepta o se rechaza la hipótesis, con los indicadores que lo demuestran. También se señalan las debilidades del estudio (Talleres 5 y 6).}

\section{Conclusión general}
\pendiente{responder el problema general e indicar si se acepta o se rechaza la hipótesis, con la reducción de costo obtenida.}

\section{Conclusiones específicas}
\begin{enumerate}[label=\textbf{\arabic*.}]
  \item \textbf{(OE1)} \pendiente{estado de la gestión de la calidad y no conformidades principales.}
  \item \textbf{(OE2)} \pendiente{línea base: índices de desperdicio y costo de no calidad de la primera losa.}
  \item \textbf{(OE3)} \pendiente{componentes del SGC diseñado.}
  \item \textbf{(OE4)} \pendiente{efecto del SGC en el costo unitario y en el desperdicio.}
  \item \textbf{(OE5)} \pendiente{resultado de las pruebas estadísticas y de la validación por expertos.}
\end{enumerate}

\section{Recomendaciones}
\guia{Recomendaciones para el Departamento de Infraestructura de la UATF, para la Supervisión y la Fiscalización, para futuras empresas contratistas y para futuras investigaciones (variables no estudiadas, otros ítems, otras obras).}

\renewcommand{\thesection}{\thechapter.\arabic{section}}
